\section{总结与延伸}

\subsection{课程全局总结表}

\begin{table}[H]
\centering
\begin{tabular}{p{0.2\textwidth} p{0.23\textwidth} p{0.24\textwidth} p{0.23\textwidth}}
\toprule
主题 & 课程给出的结论 & 关键行动 & 风险 / 关注点 \\
\midrule
企业 GenAI 趋势 & AI 民主化与平台化正在扩大用户 & 构建多模态平台、准备弹性 tooling & 依赖单一模型或 API 易造成 lock-in \\
Agent 架构 & search + LLM、Vertex AI 工具链是实作基座 & 细化 memory、retriever、planner 等模块 & 崩溃成本高、工具调用错误频发 \\
可观测与治理 & 指标、告警、审计 trail 是部署前提 & 建立 policy layer 和 human-in-loop workflow & 没有 audit trail 难以追责 \\
部署运营 & latency+cost tradeoff 需要精细调度 & multi-model routing、auto scaling 与 feedback loop & 过度扩容会浪费预算，欠扩容会影响体验 \\
人才与创新 & 需要创造力与跨学科视角 & 投资 Prompt 迭代、跨领域项目、社会价值 & 只做技术研究会错过市场与社会影响 \\
\bottomrule
\end{tabular}
\caption{Lecture 09 的核心结论与行动对照}
\end{table}

\subsection{进一步行动与反思}

\begin{itemize}
    \item 对企业产品团队：把 search 与 LLM 的组合做成一个闭环 pipeline，再嵌入质量监控与 feedback loop
    \item 对部署团队：把 tool 使用、latency、cost 做成常态仪表板，并让模型更新都在 sandbox 中先运行
    \item 对安全治理：把 Policy Layer 与 audit trail 集成，确保 agent 出错时能快速回溯
    \item 对研究者与学生：把创造力与跨学科视角融入 AI 项目，关注社会价值而非纯技术
\end{itemize}

\subsection{拓展阅读}

\begin{itemize}
    \item Google, ``Gemini: A Family of Highly Capable Multimodal Models,'' 2023.
    \item Google Cloud, ``Vertex AI Agents Documentation,'' 2024.
    \item Anthropic, ``Responsible Scaling Policy,'' 2023.
    \item Kamradt, ``Needle in a Haystack - Pressure Testing LLMs,'' 2023.
\end{itemize}

\end{document}
